\documentclass[12pt,a4paper,oneside,english]{book}

\usepackage{cite}


%\usepackage[latin1]{inputenc}

%\usepackage[T1]{fontenc}
\usepackage[english]{babel}
\usepackage{amsmath}
\usepackage{amsfonts}
\usepackage{amssymb}
\usepackage{graphicx}
\usepackage{subfig}
\usepackage{fancyhdr}
\usepackage{appendix}
\usepackage{hyphenat}
\usepackage{pdfpages}

\usepackage{array,multirow,makecell}
\newcolumntype{C}[1]{>{\arraybackslash}p{#1}}

\usepackage{enumitem}
\setlist{leftmargin=*,itemsep=0pt}

\usepackage{centernot}
\usepackage[linesnumbered,ruled,vlined,english,onelanguage]{algorithm2e}

\usepackage{quotchap}
\makeatletter
\renewcommand{\@makechapterhead}[1]{
 \chapterheadstartvskip
 {\size@chapter{\sectfont\raggedright
 {\chapnumfont
 \ifnum \c@secnumdepth >\m@ne
 \if@mainmatter\thechapter
 \fi\fi
 \par\nobreak}
 {\raggedright\advance\leftmargin10em\interlinepenalty\@M #1\par}}
 \nobreak\chapterheadendvskip}}
\makeatother
\renewcommand*{\chapterheadendvskip}{\vspace{2cm}}

\usepackage{geometry}
\geometry{hmargin=2.5cm,vmargin=2.5cm}

\pagestyle{fancyplain}
\lhead{\fancyplain{}{\nouppercase{\textit{\leftmark}}}}
\chead{\fancyplain{}{}}
\rhead{\fancyplain{}{}}
\lfoot{\fancyplain{}{}}
\cfoot{\fancyplain{}{}}
\rfoot{\fancyplain{\thepage}{\thepage}}
\renewcommand{\headrulewidth}{1pt}
\renewcommand{\footrulewidth}{1pt}

\renewcommand{\thesection}{\arabic{section}}

\usepackage{titlesec}
\titleformat{\paragraph}{\fontsize{11}{10}\bfseries}{\theparagraph}{1em}{}
\titlespacing*{\paragraph}{0pt}{10pt plus 2pt minus 0pt}{0pt plus 2pt minus 0pt}

\setcounter{secnumdepth}{4}
\setcounter{tocdepth}{4}

\usepackage{array}
\usepackage{multirow}
%\addto\captionsfrench{\def\tablename{\textsc{Tableau}}}

%\DefineBibliographyStrings{french}{urlseen = {},}

\setlength{\parskip}{8pt}
\usepackage{setspace}

\usepackage{url}

\usepackage{hyperref}
% Comment before printing to remove links' colors
\definecolor{darkblue}{rgb}{0.0, 0.0, 0.5}
\hypersetup{
 colorlinks,
 linktocpage=true,
 linkcolor={darkblue},
 citecolor={darkblue},
 urlcolor={blue}}

\sloppy

\author{You}
\title{Internship Report}

\begin{document}
\pagenumbering{gobble}
\includepdf[pages=-]{FrontPage.pdf}
\chapter*{Acknowledgments}

\frontmatter
\chapter*{Abstract}
\normalsize{Write your abstract here.

\medskip
{\noindent \textbf{Keywords: ..., ... .} }

\spacing{1}
\tableofcontents{}
\newpage 
\listoffigures
\newpage 
\listoftables
\newpage
\spacing{1.4}
\chapter*{List of acronyms}
%\addcontentsline{toc}{chapter}{Liste des acronymes}
\markboth{List of acronyms}{}
\begin{itemize}
\item \textbf{Abbrev.} Abbreviation
\end{itemize}

\mainmatter
%\mtcaddchapter[Introduction g�n�rale] 
\chapter*{Introduction}
\addcontentsline{toc}{chapter}{Introduction}
\markboth{Introduction}{}


\chapter{First chapter}
\label{ch:1er}
\section*{Introduction}

\section{Sections}
\subsection{2nd level}
\subsubsection{3rd level}
\paragraph{4th level}

\section{References}
Here's a reference to an article \cite{fisher1936use}.

Here's a reference to a Website \cite{insat}.

\section{Lists}
\subsection{Bullet lists}
Here's an example of a bullet list:
\begin{itemize}
\item Idea 1.
\item Idea 2.
\end{itemize}
\medskip \par

\subsection{Enumerated lists}
Here's an example of an enumerated list:
\begin{enumerate}
\item Idea 1 ;
\item Idea 2.
\end{enumerate}
\medskip \par

\section{Figures}
Figure \ref{fig:exemple1} is an example of a figure.
\begin{figure}[!h]
\centering
\includegraphics[width=0.3\textwidth]{images/INSAT.jpg}
\caption{This is an example of a figure.}
\label{fig:exemple1}
\end{figure}

Figure \ref{fig:exemple2} is an example of insertion of many images (\ref{fig:exemple21} and \ref{fig:exemple22}) in one figure.
\begin{figure}[!h]
\begin{center}
\subfloat[]{\includegraphics[width=0.2\textwidth]{images/INSAT.jpg}\label{fig:exemple21}}\quad 
\subfloat[]{\includegraphics[width=0.3\textwidth]{images/INSAT.jpg}\label{fig:exemple22}}
\caption{This is an example of a multi-figure: (a) small sub-figure and (b) large sub-figure.}
\label{fig:exemple2}
\end{center}
\end{figure}

\section{Tables}
Table \ref{tab:exemple1} is an example of a table with fixed columns.

\begin{table}[!h]
\caption{This is an example of a table} 
\label{tab:exemple1}
\centering
\renewcommand{\arraystretch}{1.65}
\begin{tabular}{| C{.11\textwidth} || C{.35\textwidth} | p{.45\textwidth} |} 
\hline
\textbf{a} & \textbf{b} & \textbf{c} \\ 
\hline 
\hline 
1 & \nohyphens{No hyphens avoids cutting the words.} & \nohyphens{...} \\ 
\hline 
2 & \nohyphens{...} & \nohyphens{...} \\ 
\hline 
\end{tabular}
\end{table}

Table \ref{tab:exemple2} is an example of a table where columns' widths adapt to the content.

\begin{table}[!h]
\caption{This is an example of a table with multilines et multicolumns} 
\label{tab:exemple2}
\centering
\renewcommand{\arraystretch}{1.65}
\begin{tabular}{|l||c|c|c|}
\cline{2-4}
\multicolumn{1}{c|}{\nohyphens{}} & a & b & c\\
\hline
\hline
1 & \nohyphens{...} & \nohyphens{...} & \nohyphens{...}\\ 
\hline 
2 & \multicolumn{2}{c|}{\nohyphens{...}} & \nohyphens{...}\\ 
\hline 
\multirow{2}*{3} & ... & ... & ... \\
\cline{2-4}
& \nohyphens{...} & \nohyphens{...} & \nohyphens{...}\\ 
\hline 
\end{tabular}
\end{table}

\section{Algorithms}
Algorithm \ref{algo:exemple} is an example. 
\begin{algorithm}[!h]
\setstretch{1.15}
 \caption{Process}
 \label{algo:exemple}
 \SetKwInOut{Input}{Input}
 \SetKwInOut{Output}{Output}
 \Input{$X= \{x_{i} \in\mathbb{R}^{d}\}_{i=1}^{n}$ : data matrix dimensioned $d \times n $\\
$Y = \{y_{i} \in\{1, ..., c\}\}_{i=1}^{n}$ : vector of $n$ labels\\}
\Output{$T$ : processed matrix}

... $\gets$ ...\;
\For{$i\gets 1$ \KwTo $n$ }{
\eIf{\textsf{...}}
 {
 ...\;
 }
 {
 \eIf{\textsf{\upshape ...}}
 {
 ...\;
 }
 {
 ...\;
 }
 ...
 }				 
 }
 $T \gets$ ...
\end{algorithm}

\section*{Conclusion}

\chapter*{Conclusion and perspectives}
\addcontentsline{toc}{chapter}{Conclusion and perspectives}
\markboth{Conclusion and perspectives}{}

\begin{appendix}
\chapter{Appendix 1}
Insert your appendixes here if you need.
\end{appendix}

%\spacing{1}
\bibliographystyle{unsrt}
\bibliography{references}
\end{document}